\documentclass[../HDF5_RFC.tex]{subfiles}

\begin{document}

%% ======================================================================
%%  Section 3 – Parameter String Format
%% ======================================================================
\section{Parameter String Format}
\label{sec:param-format}

\subsection{Design Rationale}

The string passed to configure a filter's parameters needs a lightweight,
unambiguous format. Several options were considered:

\begin{itemize}
\item \textbf{JSON:} Expressive but requires a parser dependency in libhdf5. Hostile in
  CLI contexts (quoting rules, no trailing commas). Heavyweight for simple
  cases.
\item \textbf{YAML:} Even heavier parsing dependency. Whitespace-significant syntax is
  problematic for single-line CLI arguments.
\item \textbf{Comma-separated key=value pairs:} Lightweight, parseable in C with no
  dependencies, consistent with how VFD/VOL configuration strings are already
  handled in HDF5 1.14.
\end{itemize}

The recommended format is comma-separated key=value pairs, with support for
bare keys (boolean flags) and double-quote-delimited values.

\subsection{Rules}
\label{sec:param-rules}

\begin{itemize}
\item Parameters are separated by commas.
\item Each parameter is a key=value pair, or a bare key (without an equals sign)
  acting as a boolean flag with an implicit empty value.
\item Leading and trailing whitespace around keys and values is stripped.
\item Keys are case-insensitive; the library normalizes to lowercase using ASCII
  case folding only (C locale \texttt{tolower()}, not locale-dependent
  \texttt{towlower()}). Values are passed to the filter callback verbatim and are
  case-sensitive by default; the filter is responsible for type interpretation.
  Note: This means \texttt{"coding=Entropy"} and \texttt{"coding=entropy"} are different
  strings from the library's perspective. Filter authors are encouraged to
  handle values case-insensitively where practical (e.g., for enum-like
  parameters) to minimize user confusion.
\item \textbf{Character encoding:} Parameter strings are restricted to printable
  ASCII (U+0021--U+007E) plus the space character (U+0020) as inter-token
  whitespace. Multi-byte encodings (UTF-8, UTF-16, etc.) are NOT supported in
  keys or unquoted values. Within quoted values, any byte sequence is accepted
  verbatim (the parser does not interpret character encoding inside quotes),
  but filter callbacks that receive non-ASCII bytes in quoted values do so at
  their own risk---the library does not guarantee encoding-safe handling. This
  restriction is a deliberate design choice to keep the parser simple,
  locale-independent, and compatible with the C-locale ASCII semantics
  described below. If future requirements demand internationalized parameter
  values, they should be passed as file-path indirection (e.g.,
  \texttt{config\_path="/path/to/config.json"}) rather than inline in the
  parameter string.
\item Keys and values consist of printable ASCII characters (excluding comma,
  semicolon, and equals sign), unless quoted.
\item \textbf{Quoted values:} If a value begins with a double quote (\texttt{"}) after
  leading whitespace has been stripped, the value extends until the next
  unescaped double quote. Within a quoted value, commas, semicolons, and
  equals signs are literal and do not act as delimiters. Whitespace inside a
  quoted value is preserved exactly as written (whitespace stripping applies
  only outside of quotes). A literal double quote within a quoted value is
  represented as two consecutive double quotes (\texttt{""}), following the
  quoted-field escaping convention of RFC~4180. An unbalanced quote
  (opening \texttt{"} with no closing \texttt{"}) MUST cause the parser to return
  \texttt{H5E\_BADVALUE} rather than silently consuming the rest of the string.
\item \textbf{Order of operations for parsing:} split the parameter string on commas
  (respecting quoted regions) $\rightarrow$ split each token on the first
  equals sign $\rightarrow$ strip leading/trailing whitespace from both key and
  value $\rightarrow$ apply quote processing to the value if it begins with a
  double quote.
\end{itemize}

\noindent\colorbox{yellow!25}{\parbox{\dimexpr\linewidth-2\fboxsep}{%
  \textbf{Open Design Decision --- Parsing Complexity:}
  Allowing commas and equals signs inside quoted values requires a stateful
  parser that tracks quoted regions before splitting on delimiters; simple
  \texttt{strtok()}-based parsing is not sufficient.
  Before finalising this design, the group should decide:
  \begin{enumerate}
    \item Do any realistic parameter values actually contain commas or equals
          signs?  If not, those characters could be \emph{prohibited} inside
          quoted values, reducing parsing to: split on \texttt{,} with
          \texttt{strtok()} $\rightarrow$ split each token on the first
          \texttt{=} $\rightarrow$ strip quotes from the value.
    \item Would a different delimiter (e.g.\ \texttt{;} or \texttt{|}) make
          quoting unnecessary entirely, given that commas would then be legal
          unquoted characters in a value such as a file path?
    \item If commas inside values are required, the current stateful-parser
          approach should be retained and an RFC~4180-compliant reference
          implementation provided to reduce implementor burden.
  \end{enumerate}
}}

\begin{itemize}
\item \textbf{Duplicate keys:} If the same key appears more than once in a parameter
  string, the parser MUST reject the string with \texttt{H5E\_BADVALUE}. The parser
  can detect duplicates cheaply during a single pass by checking each key
  against those already seen.
\end{itemize}

\noindent\colorbox{yellow!25}{\parbox{\dimexpr\linewidth-2\fboxsep}{%
  \textbf{Open Design Decision --- Duplicate Key Detection:}
  Requiring the parser to reject duplicate keys adds bookkeeping overhead
  that grows with the number of recognised keys.  While cheap today, this
  may become a concern if the set of allowed \texttt{cd\_values} expands
  significantly in the future.  The group should decide:
  \begin{enumerate}
    \item Should duplicate-key detection be \emph{removed} from the parser
          and left as a quality-of-implementation concern, accepting that
          the last (or first) occurrence silently wins?
    \item Should detection be retained but deferred to the callback layer,
          so the core parser remains a simple single-pass tokeniser with no
          per-key state?
    \item If detection is kept in the parser, should the key set be bounded
          (e.g.\ a fixed-size bitmask over known keys) to guarantee $O(1)$
          overhead regardless of input length?
  \end{enumerate}
}}

\begin{itemize}
\item \textbf{Empty keys:} A token whose key is empty after whitespace stripping MUST
  cause the parser to return \texttt{H5E\_BADVALUE}. This constraint is implicit in
  the grammar (\SectionRef{sec:param-grammar}), where \textit{key} requires
  one or more characters (\textit{printable-ascii-no-special}$^+$), but is
  stated explicitly here for implementors. The string \texttt{"=value, level=6"}
  contains an empty key before the first \texttt{=} and MUST be rejected.
\item \textbf{Key with explicit empty value vs.\ bare key:} A token of the form
  \texttt{"key="} (an equals sign present but no value after it) is a key=value
  pair with an explicit empty string value. This is syntactically distinct
  from a bare key (which has no equals sign at all). The callback receives
  an empty string \texttt{""} for the value; the bare-key convention (implicit
  empty value used as a boolean flag) applies only when no \texttt{=} is present.
  Implementations MUST preserve this distinction when dispatching to the
  callback. Example: in \texttt{"level=, fast\_mode"}, \texttt{level} has an explicit
  empty-string value and \texttt{fast\_mode} is a bare key.
\end{itemize}

Quoting is only meaningful for non-empty values; a quoted empty string
(\texttt{key=""}) MUST be rejected with \texttt{H5E\_BADVALUE}.  An empty
value is expressed as \texttt{key=} (unquoted) or as a bare key.

\begin{itemize}
\item Bare keys (without an equals sign) are permitted and treated as boolean
  flags with an implicit empty value. For example, \texttt{"fast\_mode, level=6"}
  contains two parameters: the flag \texttt{fast\_mode} (empty value) and
  \texttt{level} with value \texttt{"6"}.
\end{itemize}

\noindent\colorbox{yellow!25}{\parbox{\dimexpr\linewidth-2\fboxsep}{%
  \textbf{Open Design Decision --- Bare Keys and the Semicolon Pipeline Extension:}
  If the semicolon separator is adopted for a future pipeline-string extension
  (e.g.\ \texttt{"zstd level=3; shuffle"}), the filter name would need to be
  the first token in each semicolon-separated segment.  This means a bare key
  at position zero has a different role (filter name) from a bare key at any
  other position (boolean flag), forcing the parser to be
  position-aware---a significant complication.  The group should decide:
  \begin{enumerate}
    \item Should the pipeline extension be \emph{dropped} or deferred
          indefinitely to avoid this complexity?
    \item Should the filter name in a pipeline string be \emph{explicitly
          tagged} (e.g.\ \texttt{filter=zstd, level=3; filter=shuffle})
          so that bare keys always mean boolean flags regardless of position?
    \item Should bare keys be \emph{disallowed} entirely, requiring boolean
          flags to use an explicit value (e.g.\ \texttt{fast\_mode=1}),
          which would eliminate the position-ambiguity problem?
  \end{enumerate}
}}

\begin{itemize}
\item The semicolon restriction (outside of quoted values) is reserved for a
  possible future pipeline-string API extension.
\item A \texttt{NULL} or empty string means ``no parameters'' (valid for filters like
  shuffle and fletcher32).
\item The library defines the public constant \texttt{H5Z\_CONFIG\_STRING\_MAX}
  (= 4096) as a recommended maximum parameter string length.

\noindent\colorbox{yellow!25}{\parbox{\dimexpr\linewidth-2\fboxsep}{%
  \textbf{Note:} Implementations are not required to enforce this limit.
  Filters that require longer configuration (e.g., embedded dictionary paths)
  may exceed it; the constant is provided as a guideline for typical usage
  and for implementations that wish to bound their buffer allocation.
}}
\item The maximum number of key-value parameters in a single parameter string
  MUST NOT exceed the maximum number of \texttt{cd\_values} elements permitted
  by the library.  Strings containing more tokens than this bound MUST be
  rejected with \texttt{H5E\_BADVALUE}, as a parameter count beyond what
  \texttt{cd\_values} can represent provides no additional expressive power
  and is a potential vector for maliciously crafted strings.
\item \textbf{Locale independence:} The library's parameter string parser MUST
  be locale-independent. All character classification (whitespace detection,
  case folding) MUST use the C locale (ASCII-only \texttt{tolower()}, not
  locale-dependent \texttt{towlower()}). In particular, the parser MUST NOT be
  affected by process-wide locale settings that redefine the decimal separator
  or whitespace characters. \texttt{set\_config} callbacks that parse numeric
  values (e.g., converting a rate string to \texttt{double}) MUST likewise use
  C-locale numeric conversion (e.g., \texttt{strtod} under
  \texttt{uselocale(newlocale(LC\_NUMERIC\_MASK, "C", 0))} or a
  locale-independent equivalent such as a hand-written decimal parser). If a
  user passes \texttt{"rate=3,5"} in a comma-decimal locale, the parser splits
  on the comma and produces two tokens: \texttt{"rate=3"} and the bare key
  \texttt{"5"}. The resulting error from the filter callback may be confusing;
  documentation MUST note that parameter strings always use the period as the
  decimal separator regardless of locale.
\end{itemize}

\subsection{Examples}

\begin{center}
\begin{tabularx}{\textwidth}{l>{\raggedright\arraybackslash}X}
\toprule
\textbf{Parameter String} & \textbf{Use Case} \\
\midrule
\texttt{"level=6"} & deflate \\
\texttt{"mode=rate, rate=3.5"} & ZFP \\
\texttt{"pixels\_per\_block=16, coding=entropy"} & SZIP \\
\texttt{"scale\_type=float\_dscale, scale\_factor=4"} & scaleoffset \\
\texttt{'dict\_path="/data/run\_1,v2/dict.bin"'} & quoted value containing a comma \\
\texttt{"fast\_mode, level=6"} & bare key as boolean flag \\
\texttt{NULL} & shuffle (no params) \\
\bottomrule
\end{tabularx}
\end{center}

\subsection{Formal Grammar (EBNF)}
\label{sec:param-grammar}

The grammar below is the normative reference for the parameter string format.
Nonterminals are shown in \textit{italic}; terminals in \texttt{typewriter}
font and enclosed in single quotation marks.  Standard EBNF operators are used:
\texttt{|} for alternation, \texttt{*} for zero-or-more repetitions,
\texttt{+} for one-or-more repetitions, and parentheses for grouping.

\begin{center}
\renewcommand{\arraystretch}{1.35}
\begin{tabular}{r@{\quad$=$\quad}l}
\toprule
\textbf{Non-terminal} & \textbf{Production rule} \\
\midrule
\textit{param-string}  & \texttt{`'} $\mid$ \textit{param-list} \\
\textit{param-list}    & \textit{param} $\;(\,\text{\texttt{`,'}}\;\textit{param}\,)^{*}$ \\
\textit{param}         & \textit{key} \texttt{`='} \textit{value} $\mid$ \textit{key} \\
\textit{key}           & \textit{printable-ascii-no-special}$^{+}$ \\
\textit{value}         & \textit{bare-value} $\mid$ \textit{quoted-value} \\
\textit{bare-value}    & \textit{printable-ascii-no-special}$^{*}$ \\
\textit{quoted-value}  & \texttt{`\textquotedbl'}\,$(\,\textit{non-quote}
                         \mid \text{\texttt{`\textquotedbl\textquotedbl'}}\,)^{*}$\,\texttt{`\textquotedbl'} \\
\bottomrule
\end{tabular}
\end{center}

\noindent
where \textit{printable-ascii-no-special} is any printable ASCII character
(U+0021--U+007E) except comma~(\texttt{,}), equals~sign~(\texttt{=}),
double~quote~(\texttt{\textquotedbl}), and semicolon~(\texttt{;});
and \textit{non-quote} is any character except double~quote~(\texttt{\textquotedbl}).
An empty \textit{param-string} (the first alternative) corresponds to a
\texttt{NULL} or zero-length input, which is valid and means ``no parameters''.
Whitespace stripping (leading/trailing spaces around keys and bare values)
is applied as a post-processing step after the grammar is matched and is
not captured in the grammar itself.


%% ======================================================================
%%  Section 4 – Filter Plugin Class Version 3
%% ======================================================================
\section{Filter Plugin Class Version 3}
\label{sec:class3}

\subsection{Rationale}

The string configuration feature requires filters to participate voluntarily.
A filter must be able to accept a parameter string and produce \texttt{cd\_values},
and optionally to reverse the process for introspection. This is achieved by
introducing two new optional callbacks in a version~3 filter class structure.

\subsection{Version Numbering}
\label{sec:version-numbering}

The existing version-to-struct mapping in HDF5 is:

\begin{center}
\begin{tabular}{lll}
\toprule
\textbf{Struct} & \textbf{Version Field} & \textbf{Constant} \\
\midrule
\texttt{H5Z\_class1\_t} & none (legacy) & --- \\
\texttt{H5Z\_class2\_t} & \texttt{1} & \texttt{H5Z\_CLASS\_T\_VERS} \\
\texttt{H5Z\_class3\_t} & \texttt{2} & \texttt{H5Z\_CLASS3\_T\_VERS} \\
\bottomrule
\end{tabular}
\end{center}

The dispatch logic in \texttt{H5Zregister()} currently works by checking
\texttt{version != H5Z\_CLASS\_T\_VERS} first; any value other than
\texttt{H5Z\_CLASS\_T\_VERS} is treated as a legacy \texttt{H5Z\_class1\_t} struct.
No assumption is made about what that value will be.

This MUST be extended to also recognise \texttt{H5Z\_CLASS3\_T\_VERS} within the
\texttt{version != H5Z\_CLASS\_T\_VERS} branch, before falling back to the legacy path:

\begin{verbatim}
    if (version != H5Z_CLASS_T_VERS) {
        if (version == H5Z_CLASS3_T_VERS) -> H5Z_class3_t
        else                              -> legacy H5Z_class1_t
    } else -> H5Z_class2_t
\end{verbatim}

Note: When \texttt{H5\_NO\_DEPRECATED\_SYMBOLS} is defined, the legacy
\texttt{H5Z\_class1\_t} path is compiled out entirely and any
\texttt{version != H5Z\_CLASS\_T\_VERS} that is not \texttt{H5Z\_CLASS3\_T\_VERS}
produces a hard error. The dispatch must account for this conditional
compilation: under \texttt{H5\_NO\_DEPRECATED\_SYMBOLS}, only version~1 (v2 struct)
and \texttt{H5Z\_CLASS3\_T\_VERS} (v3 struct) are accepted.

\textbf{Strict-aliasing consideration.}
The existing dispatch in \texttt{H5Zregister()} already casts the incoming
\texttt{const void~*} to \texttt{H5Z\_class2\_t~*} to read the \texttt{version} field.
If the caller actually passed an \texttt{H5Z\_class3\_t}, this aliases one struct
type through a pointer to a different type---undefined behavior under the
C strict-aliasing rules, since the common-initial-sequence guarantee applies
only to accesses through a union, not to raw pointer casts.  This is not a new
problem introduced by v3; the same issue exists today for the v1/v2 dispatch.
However, introducing a third struct variant makes it harder to ignore.

The same issue arises in \texttt{H5PL\_\_open}, which currently casts the pointer
returned by the plugin to \texttt{H5Z\_class2\_t~*} unconditionally.  For a plugin
that returns an \texttt{H5Z\_class3\_t}, this cast is technically undefined behavior,
and the library would subsequently access fields beyond the end of what it
believes to be an \texttt{H5Z\_class2\_t}.  This RFC does not currently discuss
\texttt{H5PL\_\_open}; that omission must be addressed before the design is finalised.

Two code paths are affected and may warrant different remedies:

\begin{enumerate}
\item \textbf{Dynamic plugin loading} (\texttt{H5PL\_\_open} /
  \texttt{H5PLget\_plugin\_info}).  The library already calls
  \texttt{H5PLget\_plugin\_type()} via \texttt{dlsym} before touching the class
  struct.  Adding a parallel version-query entry point (e.g.,
  \texttt{H5PLget\_plugin\_version()}) would let the library obtain the
  interface version through a normal function call---no struct access, no
  aliasing---and cast the returned pointer to the correct type immediately.
  If no such symbol is exported, the library assumes a v1/v2 struct.

\item \textbf{Direct registration} (\texttt{H5Zregister(const void~*)}).  The
  version field must still be read from the struct.  The implementation MUST
  avoid casting to any specific \texttt{H5Z\_class} variant before the version
  is known.  Acceptable approaches include reading through a \texttt{char~*}
  cast (always permitted by the standard), \texttt{memcpy} into a local
  variable, or a compiler-specific aliasing attribute.
\end{enumerate}

The library will require internal changes to address this for both paths.
These changes are orthogonal to the v3 design but are surfaced by it.

If \texttt{H5Z\_CLASS\_T\_VERS} were bumped to~2, filter source code that
initializes \texttt{.version = H5Z\_CLASS\_T\_VERS} without guarding against version
changes would silently produce a struct the library interprets as v3, despite
missing the \texttt{set\_config}/\texttt{get\_config} fields.  However, using
\texttt{H5Z\_CLASS\_T\_VERS} as a live version token in a struct initializer is itself
incorrect practice: a filter SHOULD either hardcode the version number it was
written against, or assert at build time that \texttt{H5Z\_CLASS\_T\_VERS} has the
expected value and emit a configuration error if not.  Filter authors will need
to audit and update their code when adopting v3 regardless of which scheme is
chosen; this is an expected cost of the transition.

The current RFC therefore proposes a new constant to make the two struct
versions unambiguous without relying on filter authors having followed best
practice:

\begin{minted}{c}
#define H5Z_CLASS_T_VERS   1   /* Unchanged -- used by H5Z_class2_t */
#define H5Z_CLASS3_T_VERS  2   /* NEW -- used by H5Z_class3_t       */
\end{minted}

\noindent\colorbox{yellow!25}{\parbox{\dimexpr\linewidth-2\fboxsep}{%
  \textbf{Open Design Decision --- Struct Layout and Version Scheme:}
  Three approaches are under consideration for resolving the struct layout,
  strict-aliasing, and \texttt{H5PL\_\_open} casting issues together:
  \begin{enumerate}
    \item \textbf{Struct embedding.} Embed \texttt{H5Z\_class2\_t} as the first
          member of \texttt{H5Z\_class3\_t} and bump \texttt{H5Z\_CLASS\_T\_VERS} to~2,
          adding only the two new callback fields after the embedded base.
          A pointer to \texttt{H5Z\_class3\_t} can then be safely cast to
          \texttt{H5Z\_class2\_t~*}, resolving the strict-aliasing issue in both
          \texttt{H5Zregister()} and \texttt{H5PL\_\_open}.  However, accessing fields
          beyond the embedded base from \texttt{H5PL\_\_open} still requires knowing
          the struct version before casting to \texttt{H5Z\_class3\_t~*}.
    \item \textbf{Version-query callback.} Add an optional
          \texttt{H5PLget\_plugin\_version()} symbol to the plugin interface.
          \texttt{H5PL\_\_open} queries this via \texttt{dlsym}; if the symbol is absent,
          the plugin is assumed to export a v1/v2 struct and the existing cast
          to \texttt{H5Z\_class2\_t~*} is used.  If present, the returned version
          determines which struct type to cast to, with no aliasing violation.
          This approach avoids changes to the struct layout but adds complexity
          to the plugin loading path.
    \item \textbf{Status quo plus workaround.} Retain the current flat struct
          layout and the new \texttt{H5Z\_CLASS3\_T\_VERS} constant, and address the
          aliasing in \texttt{H5Zregister()} and \texttt{H5PL\_\_open} via
          \texttt{char~*} reads or \texttt{memcpy}, accepting that this is
          technically implementation-defined but works on all targeted platforms.
  \end{enumerate}
  Filter authors will need to update their code for any of these schemes.
  The group should decide which approach to adopt, noting that
  \texttt{H5PL\_\_open} must be explicitly covered by whatever is chosen.
}}

The version-to-struct mapping table MUST be documented in \texttt{H5Zdevelop.h}.

\verysubsection{Compatibility macro in \texttt{H5version.h}}
The existing compatibility macro
(\texttt{H5version.h}) maps the generic type alias \texttt{H5Z\_class\_t} to a versioned
struct:

\begin{minted}{c}
#if !defined(H5Z_class_t_vers) || H5Z_class_t_vers == 2
  #define H5Z_class_t H5Z_class2_t
#elif H5Z_class_t_vers == 1
  #define H5Z_class_t H5Z_class1_t
#endif
\end{minted}

Currently \texttt{H5Z\_class\_t\_vers} defaults to \texttt{2}, mapping
\texttt{H5Z\_class\_t} $\rightarrow$ \texttt{H5Z\_class2\_t}. This macro MUST NOT be changed to default to
\texttt{3} (mapping to \texttt{H5Z\_class3\_t}), because existing code using \texttt{H5Z\_class\_t}
would silently get the larger struct with uninitialized callback fields. The
default MUST remain \texttt{2}. A new value of \texttt{3} MAY be added as an opt-in for
code that explicitly requests \texttt{H5Z\_class3\_t}:

\begin{minted}{c}
#if !defined(H5Z_class_t_vers) || H5Z_class_t_vers == 2
  #define H5Z_class_t H5Z_class2_t
#elif H5Z_class_t_vers == 1
  #define H5Z_class_t H5Z_class1_t
#elif H5Z_class_t_vers == 3
  #define H5Z_class_t H5Z_class3_t
#endif
\end{minted}

Plugin authors adopting v3 SHOULD use \texttt{H5Z\_class3\_t} directly rather than
relying on the \texttt{H5Z\_class\_t} alias.


\subsection{Structure Definition}
\label{sec:class3-struct}

\begin{minted}{c}
typedef herr_t (*H5Z_set_config_func_t)(
    const char   *params,
    unsigned     *flags,
    size_t       *cd_nelmts,
    unsigned      cd_values[],
    size_t        cd_values_size
);

typedef herr_t (*H5Z_get_config_func_t)(
    unsigned          flags,
    size_t            cd_nelmts,
    const unsigned    cd_values[],
    char             *buf,
    size_t           *buf_size
);

typedef struct H5Z_class3_t {
    int                        version;
    H5Z_filter_t               id;
    unsigned                   encoder_present;
    unsigned                   decoder_present;
    const char                *name;
    const char                *description;
    H5Z_can_apply_func_t      can_apply;
    H5Z_set_local_func_t      set_local;
    H5Z_func_t                filter;
    H5Z_set_config_func_t     set_config;
    H5Z_get_config_func_t     get_config;
} H5Z_class3_t;
\end{minted}

\verysubsection{Field descriptions}

\begin{itemize}
\item \texttt{name}---The canonical identifier used for registry lookup, on-demand
  plugin discovery, and as the string accepted by
  \texttt{H5Pset\_filter\_by\_name}. MUST conform to the naming rules in
  \SectionRef{sec:param-rules}: printable ASCII, no commas, semicolons, or
  equals signs, and ideally short and lowercase (e.g., \texttt{"zfp"},
  \texttt{"blosc2"}, \texttt{"bitshuffle"}).
\item \texttt{description}---An optional human-readable string describing the
  filter (e.g., \texttt{"HDF5 ZFP filter, version 1.0.0 (rate, precision, and
  accuracy modes)"}). This field preserves the role that the \texttt{name} field
  served in \texttt{H5Z\_class2\_t}, where it was documented as a ``Comment for
  debugging'' and was returned by \texttt{H5Pget\_filter2}. For v3 plugins,
  \texttt{H5Pget\_filter2} returns the \texttt{description} field if it is
  non-\texttt{NULL}, and falls back to the \texttt{name} field otherwise. This
  ensures that applications which display the \texttt{H5Pget\_filter2} output
  continue to receive informative text after a plugin upgrades from v2 to v3.
  The \texttt{description} field may be \texttt{NULL}; no validation rules apply
  beyond requiring printable text. It is not used for registry lookup or
  any programmatic dispatch---it is purely informational.

  The dual role of the \texttt{name} field across v2 and v3 (debug comment
  in v2, canonical identifier in v3) further underscores why plugin names
  cannot be treated as authoritative identification sources without knowing
  the struct version.  Any code that consumes the \texttt{name} or
  \texttt{description} field must branch on the plugin version, and
  \texttt{H5Pget\_filter2} itself requires v2/v3 conditional logic.  This
  pattern will recur in several places in the library and potentially in
  third-party code that inspects filter metadata.  A migration guide for
  plugin authors upgrading from v2 to v3 MUST be produced alongside this
  RFC, covering at minimum: the new \texttt{name}/\texttt{description} field
  semantics, the version constant change, and the locations in the library
  where v2/v3 branching is required.
\end{itemize}

\subsection{Callback Specifications}
\label{sec:callback-specs}

\verysubsection{\texttt{H5Z\_set\_config\_func\_t}}

\begin{itemize}
\item \texttt{params}---The key-value parameter string (e.g.,
  \texttt{"mode=rate, rate=3.5"}). \texttt{NULL} if no parameters were provided.
\item \texttt{flags}---On input, contains the flags value from the API call. The
  callback MAY modify this value; the modified value is what
  \texttt{H5Pset\_filter\_by\_name} writes into the DCPL. This gives the filter the
  final say on properties like \texttt{H5Z\_FLAG\_MANDATORY} vs.\
  \texttt{H5Z\_FLAG\_OPTIONAL}. Callers that need to detect a flag override can
  compare the value they passed with the flags returned by
  \texttt{H5Pget\_filter2} after the call.
\item \texttt{cd\_nelmts}---On input, points to a value of~0. On output, set to the
  number of elements written to \texttt{cd\_values} (or the number of elements
  required, if \texttt{cd\_values} is \texttt{NULL}).
\item \texttt{cd\_values}---On output, populated with the filter's client data. If
  \texttt{NULL}, the callback MUST still set \texttt{cd\_nelmts} to the required count,
  enabling a two-pass size query pattern. The library calls the callback once
  with \texttt{cd\_values=NULL} to determine the required size, allocates
  accordingly, and calls it again with a sufficiently sized array.

  \textbf{Two-pass identity contract:} The callback MUST return the same
  \texttt{cd\_nelmts} value on both invocations. If the second-pass \texttt{cd\_nelmts}
  differs from the first-pass value in either direction, the library MUST
  return an error. A smaller second-pass value would silently truncate the
  populated array; a larger value would constitute a buffer overflow. Both
  represent a buggy callback. This contract MUST be documented in
  \texttt{H5Zdevelop.h} as part of the \texttt{H5Z\_set\_config\_func\_t} specification.

  Additionally, the library MUST validate on the second pass that
  \texttt{cd\_nelmts} does not exceed the allocated capacity; if it does, the
  library returns an error rather than writing out of bounds.

  The library MUST also verify that \texttt{cd\_nelmts} does not exceed
  \texttt{UINT16\_MAX}, the current file-format limit on the number of
  \texttt{cd\_values} elements per filter entry.  A callback returning a count
  larger than this MUST be rejected with \texttt{H5E\_BADVALUE}.  The parameter
  type is \texttt{size\_t} rather than a narrower type so that the limit can be
  raised in a future format revision without changing the callback signature.
\item \texttt{cd\_values\_size}---The capacity of the \texttt{cd\_values} array in
  bytes (0 when \texttt{cd\_values} is \texttt{NULL}).
\item Returns non-negative on success, negative on failure.
\end{itemize}

The callback SHOULD use \texttt{H5Epush} to report meaningful error messages for
invalid parameters.

\verysubsection{Unknown-key handling}
A \texttt{set\_config} callback SHOULD reject unrecognized keys by returning a
negative value and pushing \texttt{H5E\_BADVALUE} with a message identifying the
key.  This is the primary defense against misspelled parameter names (e.g.,
\texttt{"clevl=6"} instead of \texttt{"clevel=6"}), which would otherwise be
silently ignored.  Because the string API does not have the compile-time safety
of named constants, rejecting unknown keys is strongly recommended for
usability.

However, a filter MAY choose to silently ignore unrecognized keys to support
forward compatibility: a caller that targets a newer version of a plugin can
write a parameter string containing new keys and have it accepted by an older
version of the same plugin without error, falling back to defaults for the
unknown parameters.  Filter authors who choose this approach MUST document it
clearly so that callers are not misled into believing all keys were acted upon.

All built-in filter \texttt{set\_config} callbacks reject unknown keys.  This
guidance MUST be documented in \texttt{H5Zdevelop.h} as part of the
\texttt{H5Z\_set\_config\_func\_t} contract.

\verysubsection{Locale requirement for \texttt{set\_config}}
\texttt{set\_config} callbacks that parse numeric values from the parameter
string (e.g., converting a rate or tolerance to \texttt{double}) MUST use
C-locale numeric conversion. Process-wide locale settings (e.g., a
comma-decimal locale) MUST NOT affect parameter interpretation. In practice,
this means using \texttt{strtod} under \texttt{uselocale(newlocale(LC\_NUMERIC\_MASK,
"C", 0))} or a locale-independent equivalent. The library-provided helper
\texttt{H5Zconfig\_get\_param} (\SectionRef{sec:config-get-param}) returns raw
string values; it is the callback's responsibility to convert them to numeric
types in a locale-safe manner. This requirement mirrors the output-side locale
requirement on \texttt{get\_config} (\SectionRef{sec:callback-specs}) and ensures
round-trip correctness. This requirement MUST be documented in
\texttt{H5Zdevelop.h} as part of the \texttt{H5Z\_set\_config\_func\_t} specification.

\verysubsection{\texttt{H5Z\_get\_config\_func\_t}}

\begin{itemize}
\item Provides a best-effort human-readable representation of the filter
  configuration encoded in \texttt{cd\_values}.  Due to the lossy nature of the
  \texttt{cd\_values} encoding (e.g., floating-point values quantised to integers,
  information collapsed during the \texttt{set\_config} conversion), the output
  cannot in general be guaranteed to round-trip back to an identical
  configuration.  See the note on string reconstruction in
  \SectionRef{sec:introspection}.
\item Follows the standard HDF5 buffer-size query pattern: \texttt{buf\_size} is
  an in/out parameter. On entry, \texttt{*buf\_size} is the capacity of \texttt{buf};
  on return, \texttt{*buf\_size} is set to the number of bytes required (excluding
  the NUL terminator). If \texttt{buf} is \texttt{NULL} or the capacity is~0, the
  callback MUST still set \texttt{*buf\_size} to the required size and return
  success, enabling a two-pass size query.
\item Returns \texttt{herr\_t}: zero on success, negative on failure.
\item The output MUST conform to the grammar defined in
  \SectionRef{sec:param-rules} so that it is at least valid input to
  \texttt{set\_config}, even if the resulting configuration may differ from the
  original due to precision loss or encoding limitations.
\item \textbf{Locale requirement:} \texttt{get\_config} implementations that format
  numeric values (e.g., floating-point parameters) MUST use the \texttt{"C"} locale
  for number formatting. Process-wide locale settings (e.g., a locale that
  uses a comma as the decimal separator) MUST NOT affect the output. In
  practice, this means using \texttt{snprintf} with \texttt{uselocale(newlocale(LC\_NUMERIC\_MASK,
  "C", 0))} or equivalent, or formatting numeric values through a
  locale-independent method. Outputting \texttt{"rate=3,5"} in a
  comma-decimal locale would violate the parameter string grammar and break
  round-tripping via \texttt{set\_config}. This requirement MUST be documented
  in \texttt{H5Zdevelop.h} as part of the \texttt{H5Z\_get\_config\_func\_t} specification.
\end{itemize}


\subsection{Backward Compatibility}

\texttt{H5Zregister()} inspects the \texttt{version} field to determine whether the
structure is \texttt{H5Z\_class1\_t} (legacy), \texttt{H5Z\_class2\_t}, or
\texttt{H5Z\_class3\_t}. Existing v1 and v2 plugins continue to function without
modification. The new callbacks and the \texttt{description} field are never
accessed for v1/v2 registrations.
\texttt{H5Z\_class3\_t} is a new type, not a modification of \texttt{H5Z\_class2\_t}.

A v3 plugin remains fully usable through the traditional \texttt{H5Pset\_filter} API.
Upgrading to v3 is optional and additive.

For v3 plugins, the library internally stores both the \texttt{name} and
\texttt{description} fields. When \texttt{H5Pget\_filter2} is called, it returns
the \texttt{description} field in its \texttt{name[]} output parameter if
\texttt{description} is non-\texttt{NULL}, and falls back to the \texttt{name} field
otherwise. This preserves backward compatibility for applications that display
the \texttt{H5Pget\_filter2} output (\SectionRef{sec:name-semantic-change}).

If a v3 plugin's \texttt{name} field violates the naming rules of
\SectionRef{sec:param-rules}, \texttt{H5Zregister} MUST return
\texttt{H5E\_BADVALUE} and MUST NOT add the filter to the filter table.


%% ======================================================================
%%  Section 5 – Public API
%% ======================================================================
\section{Public API}
\label{sec:public-api}

\subsection{Single-Filter Configuration by Name}
\label{sec:set-filter-by-name}

\begin{minted}{c}
herr_t H5Pset_filter_by_name(
    hid_t       plist_id,
    const char *name,
    unsigned    flags,
    const char *params
);
\end{minted}

\noindent\colorbox{yellow!25}{\parbox{\dimexpr\linewidth-2\fboxsep}{%
  \textbf{Open Design Decision --- API Shape and Filter Identification:}
  The choice to identify a filter by name rather than by integer ID warrants
  scrutiny.  The application already knows which filter it wants and what
  parameters it expects, so name-based lookup does not meaningfully decouple
  the caller from the filter.  The name-registry problems discussed in
  \SectionRef{sec:name-discovery} mean that name-based identification
  introduces failure modes that ID-based identification avoids.
  Two alternatives are worth considering:
  \begin{enumerate}
    \item \textbf{\texttt{H5Pset\_filter2}}: an updated version of the existing
          function that accepts a \texttt{const char~*params} string alongside
          the existing \texttt{H5Z\_filter\_t id} and \texttt{cd\_values} arguments.
          The filter is still identified by numeric ID (no registry lookup),
          and the string is passed directly to the filter's \texttt{set\_config}
          callback.  This is the closest analogue to the existing VFD callback
          pattern and avoids all name-resolution complexity.
    \item \textbf{\texttt{H5Pset\_filter\_by\_value}}: same as above but modelled
          explicitly on the VFD \texttt{\_by\_value} naming convention, making
          the parallel intent clear.
  \end{enumerate}
  The name-based API may still be useful for scripting or configuration-file
  driven workflows where the filter ID is not known at compile time, but the
  group should decide whether that use case justifies the added complexity and
  whether it belongs in this RFC or a separate one.
}}

\verysubsection{Description}
Looks up the filter identified by \texttt{name}, configures it using
the key-value parameter string \texttt{params}, and appends it to the existing filter
pipeline on \texttt{plist\_id}. If a filter with the same ID is already present in the
pipeline, it is replaced (not duplicated)---this matches the existing
behavior of \texttt{H5Pset\_filter}. The pipeline is subject to the existing maximum
of \texttt{H5Z\_MAX\_NFILTERS} (32) entries.

This function is modeled on \texttt{H5Pset\_filter}: it appends a single filter to
the pipeline with explicit flag control. The difference is that the filter is
identified by name rather than integer ID, and parameters are passed as a
human-readable string rather than a \texttt{cd\_values} array. The two APIs are fully
interchangeable on the same DCPL---calling \texttt{H5Pset\_filter\_by\_name} and
\texttt{H5Pset\_filter} for the same filter ID produces identical results in the
property list, and either can replace a filter entry created by the other.

\verysubsection{Behavioral note on flags}
Unlike \texttt{H5Pset\_filter}, where the caller has
absolute control over flags, \texttt{H5Pset\_filter\_by\_name} passes the caller's
flags through the filter's \texttt{set\_config} callback, which MAY modify them. This
allows filters to enforce constraints (e.g., SZIP unconditionally forces
\texttt{H5Z\_FLAG\_OPTIONAL}). This is a deliberate asymmetry: filters configured via
string may override the caller's flag choice based on internal requirements.
Callers that need to detect a flag override can compare the flags they passed
with the flags returned by \texttt{H5Pget\_filter2} after the call. See
\SectionRef{sec:callback-specs}.

If the filter name is not found in the in-memory name registry, the library
triggers an on-demand plugin search (see \SectionRef{sec:name-discovery}).

\verysubsection{Parameters}

\begin{itemize}
\item \texttt{plist\_id}---A dataset creation property list identifier.
\item \texttt{name}---A filter name (case-insensitive, e.g., \texttt{"deflate"},
  \texttt{"zfp"}) or a string representation of a numeric filter ID (e.g.,
  \texttt{"32013"}). Resolution order: the name registry is checked first; if the
  name is not found and consists entirely of digits, it is parsed as a numeric
  filter ID. This means a registered name takes precedence even if it happens
  to be a numeric string.
\item \texttt{flags}---A bit mask specifying filter behavior
  (\texttt{H5Z\_FLAG\_MANDATORY}, \texttt{H5Z\_FLAG\_OPTIONAL}).
\item \texttt{params}---A comma-separated key=value parameter string, or \texttt{NULL}
  for filters that take no parameters. If \texttt{params} is non-\texttt{NULL} and
  non-empty for a filter that takes no parameters, the filter's
  \texttt{set\_config} callback SHOULD reject it with \texttt{H5E\_BADVALUE}. This is the
  expected pattern for parameterless built-in filters (shuffle, fletcher32,
  nbit).
\end{itemize}

\verysubsection{Returns}
Non-negative on success; negative on failure.

\verysubsection{Errors}

\begin{itemize}
\item \texttt{H5E\_NOFILTER}---The filter name could not be resolved to a registered
  filter, even after an on-demand plugin search.
\item \texttt{H5E\_UNSUPPORTED}---The filter was found but does not implement
  \texttt{set\_config} and \texttt{params} is not \texttt{NULL}. The user should use
  \texttt{H5Pset\_filter} with raw \texttt{cd\_values} instead.
\item \texttt{H5E\_BADVALUE}---The parameter string is malformed or rejected by the
  filter's \texttt{set\_config} callback.
\item \texttt{H5E\_CANTINIT}---The \texttt{set\_config} callback returned more
  \texttt{cd\_values} elements than the library can accommodate.
  Note: \texttt{H5Z\_COMMON\_CD\_VALUES} is~4, not a hard cap; the current code
  permits this to be exceeded regardless of dataset layout.  The precise
  upper bound and the conditions under which it applies need to be
  verified against the current implementation before this error condition
  can be specified normatively.
\end{itemize}

\verysubsection{Usage example}

\begin{minted}{c}
// Current approach:
H5Pset_filter(plist, H5Z_FILTER_DEFLATE, H5Z_FLAG_MANDATORY, 1, &level);

// Proposed equivalent:
H5Pset_filter_by_name(plist, "deflate", H5Z_FLAG_MANDATORY, "level=6");

// Plugin filter -- no compiled-in constants needed:
H5Pset_filter_by_name(plist, "zfp", H5Z_FLAG_MANDATORY,
                       "mode=rate, rate=3.5");

// Building a pipeline with multiple calls:
H5Pset_filter_by_name(plist, "shuffle", H5Z_FLAG_MANDATORY, NULL);
H5Pset_filter_by_name(plist, "deflate", H5Z_FLAG_MANDATORY, "level=9");
H5Pset_filter_by_name(plist, "fletcher32", H5Z_FLAG_MANDATORY, NULL);
\end{minted}


\subsection{Pipeline Introspection}
\label{sec:introspection}

\begin{minted}{c}
herr_t H5Pget_filter_name_by_idx(
    hid_t     plist_id,
    unsigned  filter_idx,
    char     *name_buf,
    size_t    name_buf_size,
    size_t   *name_len,
    unsigned *flags
);

herr_t H5Pget_filter_params_by_idx(
    hid_t     plist_id,
    unsigned  filter_idx,
    char     *params_buf,
    size_t    params_buf_size,
    size_t   *params_len
);
\end{minted}

\verysubsection{Description}
These two functions retrieve the string representation of a
single filter in the pipeline at position \texttt{filter\_idx}.

\texttt{H5Pget\_filter\_name\_by\_idx} returns the filter's canonical name and flags. If
the filter is registered in the name registry, its canonical name is written
to \texttt{name\_buf}. Otherwise, the numeric filter ID is formatted as a decimal
string (e.g., \texttt{"32013"}).

\texttt{H5Pget\_filter\_params\_by\_idx} returns the human-readable parameter string.
For filters that implement \texttt{get\_config}, the callback's output is used. For
filters that do not implement \texttt{get\_config}, the raw \texttt{cd\_values} are formatted
as a fallback string using colon separators (e.g.,
\texttt{"cd\_values=1:0:0:1074921472"}). The colon separator avoids ambiguity with
the comma-separated key=value grammar defined in \SectionRef{sec:param-format}---a fallback
string is a single key=value pair where the value contains colons. This means
the fallback string is a valid parameter string under the grammar, preserving
the possibility of round-tripping in a future release.

\verysubsection{Note on string reconstruction}
The DCPL does not retain the original
parameter string. After \texttt{H5Pset\_filter\_by\_name} resolves the string into
\texttt{cd\_values}, only the \texttt{cd\_values} are stored in the property list (identical
to \texttt{H5Pset\_filter}). These functions rely on the filter's \texttt{get\_config}
callback to reconstruct a human-readable string from the stored \texttt{cd\_values}.

This reconstruction is inherently limited. The \texttt{cd\_values} encoding was
designed to carry opaque integer data, not to serve as a lossless serialisation
of a key-value string. For some filters (e.g., ZFP), floating-point
configuration values are quantised during conversion to unsigned integers; the
original floating-point value cannot be exactly recovered from the stored
integer, and the reconstructed string will reflect the rounded value rather
than what the user originally typed. Consumers MUST NOT rely on exact string
equality between input and output, and SHOULD treat the reconstructed string
as informational only.

\noindent\colorbox{yellow!25}{\parbox{\dimexpr\linewidth-2\fboxsep}{%
  \textbf{Open Design Decision --- On-disk storage of configuration:}
  True round-tripping requires changing what is stored on-disk. Three
  approaches are worth considering:
  \begin{enumerate}
    \item \textbf{Store the raw parameter string.} Simple to implement, but
          potentially wasteful of space for long strings, and displays the
          user's original input rather than the post-processing values
          (e.g., a rounded rate) that the filter actually used.
    \item \textbf{Store tokenized key-value pairs.} The library tokenizes the
          string once and stores the key-value pairs directly in the pipeline
          message. More compact than the raw string, and consistent with the
          proposal that tokenization should be the library's responsibility
          rather than each filter's. The stored pairs would reflect the
          original input values, not the post-processing values.
    \item \textbf{Store post-processing key-value pairs.} After the filter's
          \texttt{set\_config} callback runs, the filter reports back the
          effective key-value pairs (e.g., the rounded rate that was actually
          applied). These are stored in the pipeline message. Tools like
          \texttt{h5dump} would display what the filter actually used, not what
          the user typed---arguably more useful, but requires a new callback
          contract for filters to report their effective configuration.
  \end{enumerate}
  All three options require a format change. The group should decide whether
  any of them is in scope for this RFC or should be deferred to a follow-on.
}}

\verysubsection{Note on \texttt{H5Pget\_filter2}}
The existing \texttt{H5Pget\_filter2} function returns a
filter name from the class struct's \texttt{name} field (the ``debug comment'' in v2
plugins). For v3 plugins, \texttt{H5Pget\_filter2} returns the \texttt{description}
field if it is non-\texttt{NULL}, and falls back to the \texttt{name} field otherwise.
This preserves the existing behavior of returning human-readable descriptive
text.

\texttt{H5Pget\_filter\_name\_by\_idx} returns the canonical name from the name
registry, which may differ from the \texttt{H5Pget\_filter2} output. For example, a
v3 plugin might have \texttt{name = "zfp"} and
\texttt{description = "HDF5 ZFP filter, version 1.0.0 (rate mode)"} in its class
struct; \texttt{H5Pget\_filter2} returns the description, while
\texttt{H5Pget\_filter\_name\_by\_idx} returns \texttt{"zfp"}. The two functions serve
different purposes and may return different strings for the same pipeline entry.

\noindent\colorbox{yellow!25}{\parbox{\dimexpr\linewidth-2\fboxsep}{%
  \textbf{Open Design Decision --- \texttt{H5Pget\_filter2} semantics and new introspection API:}
  Changing \texttt{H5Pget\_filter2} to return \texttt{description} (falling back to
  \texttt{name}) alters the semantics of an existing function in a way that may
  surprise callers, and the benefit is unclear given that the
  \texttt{name}/\texttt{description} split is itself a consequence of the v3 design
  choices.  Similarly, \texttt{H5Pget\_filter\_name\_by\_idx} is only meaningful if
  there is a concrete use case for retrieving the canonical registered name
  independently of the human-readable description.  The group should decide:
  \begin{enumerate}
    \item Should \texttt{H5Pget\_filter2} behaviour be left unchanged for v3
          plugins (always returning \texttt{name}), avoiding any semantic drift?
    \item Is there a demonstrated use case for \texttt{H5Pget\_filter\_name\_by\_idx}
          that justifies adding it, or should it be deferred until a concrete
          need is identified?
    \item If the name-based API (\SectionRef{sec:set-filter-by-name}) is
          reconsidered in favour of an ID-based approach, both of these
          introspection functions may become unnecessary.
  \end{enumerate}
}}

\verysubsection{Parameters (\texttt{H5Pget\_filter\_name\_by\_idx})}

\begin{itemize}
\item \texttt{plist\_id}---A dataset creation property list identifier.
\item \texttt{filter\_idx}---Zero-based index of the filter in the pipeline
  (\texttt{unsigned}, consistent with \texttt{H5Pget\_filter2}).
\item \texttt{name\_buf}---Buffer to receive the filter's canonical name, or
  \texttt{NULL} to query size only.
\item \texttt{name\_buf\_size}---Size of \texttt{name\_buf} in bytes.
\item \texttt{name\_len}---On output, set to the length of the name string
  (excluding NUL), or \texttt{NULL} if not needed. Set even when \texttt{name\_buf}
  is \texttt{NULL}, enabling a size query.
\item \texttt{flags}---Pointer to receive the filter's flags, or \texttt{NULL}.
\end{itemize}

\verysubsection{Parameters (\texttt{H5Pget\_filter\_params\_by\_idx})}

\begin{itemize}
\item \texttt{plist\_id}---A dataset creation property list identifier.
\item \texttt{filter\_idx}---Zero-based index of the filter in the pipeline.
\item \texttt{params\_buf}---Buffer to receive the parameter string, or \texttt{NULL}
  to query size only.
\item \texttt{params\_buf\_size}---Size of \texttt{params\_buf} in bytes.
\item \texttt{params\_len}---On output, set to the length of the parameter string
  (excluding NUL), or \texttt{NULL} if not needed. Set even when \texttt{params\_buf}
  is \texttt{NULL}, enabling a size query.
\end{itemize}

\verysubsection{Returns (both functions)}
Non-negative on success; negative on failure.

These functions complement the existing \texttt{H5Pget\_filter2} by providing the same
information in human-readable form.


\subsection{Developer Helper: Parameter Lookup}
\label{sec:config-get-param}

\begin{minted}{c}
htri_t H5Zconfig_get_param(
    const char *params,
    const char *key,
    char       *value_buf,
    size_t     *buf_size
);
\end{minted}

\verysubsection{Description}
Searches the parameter string \texttt{params} for the given \texttt{key}
and returns the associated value. This function is declared in
\texttt{H5Zdevelop.h} and is the library-provided parameter parser that plugin
authors are expected to use in their \texttt{set\_config} callbacks. All built-in
filter \texttt{set\_config} implementations (deflate, szip, scaleoffset) use this
function.

Returns positive (\texttt{> 0}) if the key is present in \texttt{params}, \texttt{0} if the key
is not present, or negative on error (malformed string). The \texttt{htri\_t} return
type (HDF5's standard tri-state boolean) cleanly separates three cases: found,
not found, and error---without triggering the error stack for a missing
optional parameter.

\texttt{buf\_size} is an in/out parameter following the standard HDF5 size-query
pattern. On entry, \texttt{*buf\_size} is the capacity of \texttt{value\_buf}. On
return, \texttt{*buf\_size} is set to the number of bytes required for the processed
value (excluding NUL), allowing callers to allocate correctly on a first pass.
If \texttt{value\_buf} is \texttt{NULL} or \texttt{*buf\_size} is~0 on entry, the function
performs a size query only and does not write to \texttt{value\_buf}.  When the key
is found and \texttt{value\_buf} is non-\texttt{NULL}, the processed value is written as a
NUL-terminated string (quotes stripped, \texttt{""} unescaped).  The caller can use
\texttt{strlen} on the result; no separate length output is provided.

A bare key like \texttt{"fast\_mode"} (no equals sign) returns positive with
\texttt{*buf\_size = 0} and an empty string in \texttt{value\_buf}, distinguishable
from an absent key (returns \texttt{0}).

This function is placed in \texttt{H5Zdevelop.h} (developer-facing) rather than
\texttt{H5Zpublic.h} to limit the public API surface.


\subsection{Canonical \texttt{set\_config} Implementation Template}
\label{sec:set-config-template}

The two-pass \texttt{set\_config} protocol (\SectionRef{sec:callback-specs}) is
subtle. Plugin authors MUST ensure identical \texttt{cd\_nelmts} on both passes.
The following template shows the recommended implementation pattern. All
built-in filter \texttt{set\_config} callbacks follow this structure and are
documented in \texttt{H5Zdevelop.h} as normative examples.

\begin{minted}{c}
static herr_t
my_filter_set_config(const char *params, unsigned *flags,
                     size_t *cd_nelmts, unsigned cd_values[],
                     size_t cd_values_size)
{
    /* ---- 1. Parse parameters (identical on both passes) ---- */
    char   val_buf[64];
    size_t buf_size = sizeof(val_buf);
    int    level  = 6;         /* default */

    if (params) {
        htri_t found = H5Zconfig_get_param(params, "level",
                                           val_buf, &buf_size);
        if (found < 0) return FAIL;
        if (found > 0) {
            char *end;
            long v = strtol(val_buf, &end, 10);
            if (*end != '\0' || v < 0 || v > 9)
                return FAIL;           /* push H5E_BADVALUE */
            level = (int)v;
        }
    }

    /* ---- 2. Report element count (SAME on both passes) ---- */
    *cd_nelmts = 1;

    /* ---- 3. Populate cd_values only on the second pass ---- */
    if (cd_values != NULL) {
        if (cd_values_size < *cd_nelmts * sizeof(*cd_values))
            return FAIL;               /* should not happen */
        cd_values[0] = (unsigned)level;
    }

    /* ---- 4. Optionally override flags ---- */
    /* (void)flags;  -- accept caller's flags unchanged */

    return SUCCEED;
}
\end{minted}

\verysubsection{Key invariants}
\begin{itemize}
\item \texttt{cd\_nelmts} MUST be set to the same value regardless of whether
  \texttt{cd\_values} is \texttt{NULL} or non-\texttt{NULL}.
\item Parameter parsing MUST be performed on both passes (not cached) to
  ensure the count is consistent.
\item If \texttt{cd\_values} is non-\texttt{NULL}, verify
  \texttt{cd\_values\_size >= *cd\_nelmts * sizeof(*cd\_values)} before writing.
\item For floating-point conversions, use C-locale \texttt{strtod} (see locale
  requirements in \SectionRef{sec:param-rules} and \SectionRef{sec:callback-specs}).
\end{itemize}

\textbf{Note on double parsing:} Because \texttt{cd\_nelmts} may vary depending on
the parameter values (not just their presence), the two-pass protocol requires
the callback to parse the parameter string twice.  Filters that wish to avoid
this overhead may allocate a worst-case \texttt{cd\_values} buffer on the stack
(bounded by \texttt{UINT16\_MAX} but in practice much smaller) and skip the
size-query pass by always reporting the maximum count, filling unused slots
with a sentinel value.  The library does not mandate which approach is used;
both are conforming.


\subsection{Design Alternative: Pipeline-in-a-String API}
\label{sec:pipeline-string}

An alternative (or additional) API has been proposed in community discussion
that allows an entire filter pipeline to be expressed as a single string:

\begin{minted}{c}
H5Pset_filter_pipeline(plist, "shuffle; deflate, level=9; fletcher32");
\end{minted}

In this variant, semicolons separate individual filter specifications, and the
first token in each specification is the filter name. The function would
replace the entire pipeline atomically.

\verysubsection{Arguments in favor}

\begin{itemize}
\item Very concise for CLI tools and configuration files.
\item A single string can be stored, logged, and passed around trivially.
\item Mirrors how users think about a pipeline as a unit.
\end{itemize}

\verysubsection{Arguments against}

\begin{itemize}
\item Introduces a grammar (semicolons, ordering rules, escaping) that the
  library must specify, parse, and maintain indefinitely.
\item Departs from the established HDF5 pattern of one-filter-per-call with
  explicit parameters per call.
\item Makes per-filter flag control awkward (all filters share a default, or an
  inline syntax must be invented).
\item Conflates filter selection and parameter passing in a way that is harder to
  validate and produces less precise error messages.
\end{itemize}

\verysubsection{Recommendation}
This document proposes \texttt{H5Pset\_filter\_by\_name} (\SectionRef{sec:set-filter-by-name})
as the primary API, following the established HDF5 pattern. The
pipeline-string variant could be added as a convenience function in a future
release if demand warrants it, since it can be trivially implemented as a
wrapper that tokenizes on semicolons and calls \texttt{H5Pset\_filter\_by\_name} in a
loop.

\verysubsection{Note on grammar interaction}
The bare-key feature (\SectionRef{sec:param-rules}) creates an
ambiguity for the pipeline-string grammar. In a pipeline string like
\texttt{"shuffle; deflate, fast\_mode, level=9"}, it is unclear whether
\texttt{"fast\_mode"} is a second filter name or a bare-key parameter of deflate. If
the pipeline-string variant is added in the future, the grammar MUST resolve
this---the natural rule is that the filter name is always the first token
after a semicolon delimiter, and all subsequent comma-separated tokens in the
same spec are parameters.


%% ======================================================================
%%  Section 6 – Filter Name Registry
%% ======================================================================
\section{Filter Name Registry}
\label{sec:name-registry}

\subsection{Overview}

The library maintains an internal mapping from canonical names to filter IDs
(\texttt{H5Z\_filter\_t}). This registry is the mechanism by which the new APIs resolve
filter names at runtime. It also serves as a cache for names discovered during
plugin loading.

\subsection{Built-in Registrations}

The following names are registered at library initialization:

\begin{center}
\begin{tabular}{ll}
\toprule
\textbf{Name} & \textbf{Filter ID} \\
\midrule
\texttt{"deflate"}     & \texttt{H5Z\_FILTER\_DEFLATE} \\
\texttt{"shuffle"}     & \texttt{H5Z\_FILTER\_SHUFFLE} \\
\texttt{"fletcher32"}  & \texttt{H5Z\_FILTER\_FLETCHER32} \\
\texttt{"szip"}        & \texttt{H5Z\_FILTER\_SZIP} \\
\texttt{"nbit"}        & \texttt{H5Z\_FILTER\_NBIT} \\
\texttt{"scaleoffset"} & \texttt{H5Z\_FILTER\_SCALEOFFSET} \\
\bottomrule
\end{tabular}
\end{center}

\subsection{Plugin Auto-Registration}

When a filter plugin compiled against \texttt{H5Z\_class3\_t} is loaded (whether by
integer-ID search or name-based search), the library reads the \texttt{name} field
from the class structure and automatically registers it in the name registry.
If the \texttt{name} field is \texttt{NULL}, the plugin is not registered by name and can
only be addressed by numeric ID.

The library MUST make an internal copy (\texttt{strdup}) of the name string when
inserting it into the name registry. The \texttt{name} field in \texttt{H5Z\_class3\_t} points
to memory inside the loaded shared library. If the plugin is later unloaded
(e.g., via \texttt{H5Zunregister} or \texttt{H5close}), the registry would otherwise
contain a dangling pointer.

\verysubsection{Note on the \texttt{name} and \texttt{description} fields}
In \texttt{H5Z\_class2\_t}, the \texttt{name} field is
documented as ``Comment for debugging'' but is also user-visible---it is the
string returned by \texttt{H5Pget\_filter2} in its \texttt{name[]} output parameter.
Existing plugins may contain freeform text (e.g., ``my cool zfp filter v2.1'').
In \texttt{H5Z\_class3\_t}, the \texttt{name} field takes on a stronger role as a canonical
identifier used for registry lookup and on-demand plugin discovery, and MUST
follow the naming rules in \SectionRef{sec:param-rules} (printable ASCII, no
commas/semicolons/equals signs). Plugin authors upgrading from v2 to v3 MUST
choose a suitable canonical name. This semantic shift MUST be prominently
documented in the migration guide.

The new \texttt{description} field in \texttt{H5Z\_class3\_t} preserves the display
role that the \texttt{name} field served in v2. Plugin authors upgrading from v2
to v3 SHOULD move their existing freeform \texttt{name} text (e.g.,
\texttt{"HDF5 ZFP filter, version 1.0.0 (rate mode)"}) into the
\texttt{description} field, and assign a short canonical identifier to
\texttt{name} (e.g., \texttt{"zfp"}). This ensures continuity for applications
that display the output of \texttt{H5Pget\_filter2}
(\SectionRef{sec:name-semantic-change}).


\subsection{Explicit Registration and Lookup}
\label{sec:explicit-registration}

\noindent\colorbox{yellow!25}{\parbox{\dimexpr\linewidth-2\fboxsep}{%
  \textbf{Open Design Decision --- Application Access to the Name Registry:}
  Exposing \texttt{H5Zregister\_name} to applications gives them the ability to
  remap arbitrary names to arbitrary filter IDs (e.g., mapping \texttt{"zfp"} to
  the ID for zlib), creating opportunities for misconfiguration that would be
  difficult to diagnose.  It also complicates support, since the registry
  state becomes part of the application's environment rather than a
  library-controlled invariant.  Unless a concrete use case is presented that
  cannot be served by implicit self-registration (via \texttt{H5Zregister} with an
  \texttt{H5Z\_class3\_t}) or on-demand plugin loading, the group should consider:
  \begin{enumerate}
    \item \textbf{Removing \texttt{H5Zregister\_name} entirely} from the public API,
          leaving name registration as an internal operation performed only
          by \texttt{H5Zregister} when it processes an \texttt{H5Z\_class3\_t}.
    \item \textbf{Restricting it to \texttt{H5Zdevelop.h}} so it is available to
          plugin authors for testing but not to general application code.
    \item \textbf{Retaining it} only if a specific use case (e.g., registering
          aliases for v2 plugins that cannot self-register) justifies the risk.
  \end{enumerate}
  \texttt{H5Zlookup\_by\_name} and \texttt{H5Zfilter\_avail\_by\_name} are read-only and
  do not carry the same risk; the concern applies specifically to the
  write path (\texttt{H5Zregister\_name}).
}}

\begin{minted}{c}
herr_t H5Zregister_name(H5Z_filter_t id, const char *canonical_name);
\end{minted}

Registers a name mapping. This can be used to add aliases (e.g.,
\texttt{"gzip"} $\rightarrow$ \texttt{H5Z\_FILTER\_DEFLATE}) or to register names for older v2 plugins
that cannot self-register. A single filter ID MAY be associated with multiple
names. In this case, the introspection functions (\SectionRef{sec:introspection}) return the name
that was first registered for that ID (either the built-in name, the name from
the plugin's class struct, or the first explicit registration).
Returns non-negative on success or negative on error (e.g., \texttt{id} is
not currently registered, or \texttt{canonical\_name} violates the naming rules
in \SectionRef{sec:param-rules}).

\begin{minted}{c}
H5Z_filter_t H5Zlookup_by_name(const char *name);
\end{minted}

Returns the filter ID associated with \texttt{name}, or \texttt{H5Z\_FILTER\_ERROR} if not
found. Lookup is case-insensitive (ASCII folding only). This function queries
the in-memory name registry only. It does NOT trigger on-demand plugin
loading. It is a fast, pure lookup with no I/O side effects.

\begin{minted}{c}
htri_t H5Zfilter_avail_by_name(const char *name);
\end{minted}

Returns positive if the filter identified by \texttt{name} is available (registered
or discoverable via plugin search), \texttt{0} if not available, negative on error.
Unlike \texttt{H5Zlookup\_by\_name}, this function DOES trigger on-demand plugin
loading if the name is not already in the registry---it is the by-name
equivalent of the existing \texttt{H5Zfilter\_avail()}. This allows users to test
filter availability before committing to a DCPL configuration.

\verysubsection{Registration Ordering Constraint and Pre-Registration}
\label{sec:preregister-name}

\texttt{H5Zregister\_name} requires that the filter identified by \texttt{id} already be
present in the filter table at the time of the call. The required call
sequence for a filter that must be explicitly named is:

\begin{enumerate}
\item Register the filter class: \texttt{H5Zregister(\&my\_class3)} (or allow
  auto-registration during plugin load, which populates the \texttt{name} field
  automatically for v3 plugins).
\item Register the name alias: \texttt{H5Zregister\_name(MY\_FILTER\_ID, "myalias")}.
\end{enumerate}

Calling \texttt{H5Zregister\_name} before \texttt{H5Zregister} MUST fail with
\texttt{H5E\_NOTFOUND}. This ordering requirement MUST be documented in the API
reference for \texttt{H5Zregister\_name}.

In some initialization sequences, the application knows the name-to-ID
mapping before the filter plugin is loaded. A deferred-binding variant is
therefore introduced to support these cases.

\verysubsection{Motivating use case}
A scientific facility operates a site-wide HDF5 configuration file
(\texttt{hdf5.conf}) that maps canonical names to filter IDs for all
facility-approved compression plugins. The facility's initialization library
reads this file at startup and pre-registers all mappings before any HDF5 file
is opened. Plugins are loaded on demand later, when a dataset actually
requests a particular filter. Without pre-registration, the initialization
library would need to either (a)~force-load all plugin shared libraries at
startup (expensive and fragile on systems with hundreds of plugins), or
(b)~defer all name registration until the first use of each plugin (losing the
ability to validate the configuration file at startup). Pre-registration
allows the facility to validate name-to-ID mappings eagerly while keeping
plugin loading lazy.

A second use case arises in test harnesses that need to register a name alias
for a mock filter before the mock is registered, so that the code under test
can use \texttt{H5Pset\_filter\_by\_name} with the expected name.

\noindent\colorbox{yellow!25}{\parbox{\dimexpr\linewidth-2\fboxsep}{%
  \textbf{Open Design Decision --- Validity of these use cases:}
  Both use cases above have been questioned.  For the facility use case, the
  name-to-ID mapping could be maintained entirely on the facility's side:
  the initialization library already reads \texttt{hdf5.conf} and holds the
  mapping in memory, so it can resolve a name to an ID itself and pass the
  ID directly to HDF5 at the point of use.  Storing the mapping inside HDF5
  does not change when plugins are loaded---lazy loading is orthogonal to
  where the name-to-ID table lives.  The validation-at-startup benefit can be
  achieved by the facility library validating its own in-memory table without
  involving HDF5 at all.  For the test-harness use case, this amounts to
  testing the pre-registration API itself rather than providing a genuine
  independent motivation.  The group should decide whether these use cases
  justify the added API surface and registry complexity, or whether
  \texttt{H5Zpreregister\_name} (and \texttt{H5Zregister\_name}) should be removed.
}}

\begin{minted}{c}
herr_t H5Zpreregister_name(H5Z_filter_t id, const char *canonical_name);
\end{minted}

\texttt{H5Zpreregister\_name} records a pending name-to-ID mapping in a separate
pre-registration table. Unlike \texttt{H5Zregister\_name}, it does NOT require
\texttt{id} to be currently registered in the filter table. The name is validated
against the naming rules of \SectionRef{sec:param-rules} at pre-registration
time; a name that violates those rules causes \texttt{H5E\_BADVALUE} to be returned
immediately, before any filter registration has occurred.

\textbf{Activation:} When \texttt{H5Zregister()} completes successfully for a filter
with a given \texttt{id}, the library MUST scan the pre-registration table and
promote any pending entries whose \texttt{id} matches. Each matching entry is moved
from the pre-registration table into the live name registry, subject to the
normal collision policy (\SectionRef{sec:collision-policy}). Promotion is
performed as an atomic step at the end of a successful \texttt{H5Zregister} call,
before that call returns.

\textbf{Ordering and priority:} A pre-registered name that is promoted at
\texttt{H5Zregister} time is treated as having been registered immediately after the
\texttt{H5Zregister} call. It therefore arrives after any name that was
auto-registered from the v3 plugin's own \texttt{name} field. Under the
first-registered-wins collision policy (\SectionRef{sec:collision-policy}),
the plugin's own \texttt{name} field takes precedence over a pre-registered alias
for the same name. If the application needs a specific alias to be the
primary (introspection-visible) name, it must ensure the alias does not
collide with the plugin's own \texttt{name} field.

\textbf{Cleanup:} Entries remaining in the pre-registration table when
\texttt{H5close()} is called (i.e., names pre-registered for filters that were
never registered) are silently discarded. Applications that depend on a
pre-registered name being available MUST verify availability via
\texttt{H5Zfilter\_avail\_by\_name} or \texttt{H5Zlookup\_by\_name} before use.
\emph{Implementor note: the internal table entries are freed via
\texttt{H5MM\_xfree} during \texttt{H5close} teardown; ``silently discarded''
refers to the absence of error reporting, not to a memory leak.}

\textbf{Collision behavior:} If a promoted name already exists in the live
registry at promotion time, the promotion is silently discarded (first-wins,
consistent with \SectionRef{sec:collision-policy}). Duplicate pre-registrations
for the same (id, name) pair are likewise silently discarded.

\textbf{Summary of call sequences:}

\begin{center}
\begin{tabularx}{\textwidth}{llX}
\toprule
\textbf{Scenario} & \textbf{API} & \textbf{Constraint} \\
\midrule
Alias after registration & \texttt{H5Zregister\_name} & Filter must be registered first \\
Alias before registration & \texttt{H5Zpreregister\_name} & Filter may be registered later \\
v3 plugin self-names & Auto-registration & Happens inside \texttt{H5Zregister} \\
\bottomrule
\end{tabularx}
\end{center}


\subsection{Collision Policy}
\label{sec:collision-policy}

If a name is already registered, \texttt{H5Zregister\_name} and plugin
auto-registration MUST fail and return an error. First registration wins. This
prevents silent behavioral changes when plugins are loaded in different orders.

Note: The collision policy applies to names, not to IDs. Multiple names may
map to the same filter ID (aliases). But a given name may map to only one ID.

\subsubsection{Interaction with H5Zunregister}

When \texttt{H5Zunregister()} removes a filter, all name registry entries pointing to
that filter's ID MUST also be removed. If the entries persist,
\texttt{H5Zlookup\_by\_name} would return a filter ID that is no longer registered, and
a subsequent \texttt{H5Pset\_filter} with that ID would fail---producing a confusing
error path. The reverse mapping (ID $\rightarrow$ names) required for this cleanup is an
implementation requirement; see \SectionRef{sec:registry-impl}. The cleanup routine MUST free all
internally copied name strings (via \texttt{H5MM\_xfree} or equivalent) to prevent
memory leaks in long-running processes or test suites that repeatedly register
and unregister filters.


\subsection{Name-Based Plugin Discovery}
\label{sec:name-discovery}

This section describes how \texttt{H5Pset\_filter\_by\_name} finds and loads a plugin
when the requested name is not yet in the registry. The mechanism extends the
existing VFD/VOL by-name plugin search pattern.

The VFD and VOL subsystems already solve this exact problem using a
discriminated union in \texttt{H5PL\_key\_t}, a \texttt{kind} field that selects between
by-name and by-value matching, and a \texttt{check\_plugin\_load} callback that
performs the match. The filter case simply has not been upgraded yet. This
design applies the same pattern.

\verysubsection{Current state}
The filter key in \texttt{H5PL\_key\_t} (\texttt{H5PLprivate.h}) is a bare
\texttt{int}:

\begin{minted}{c}
typedef union H5PL_key_t {
    int            id;            /* Filters: integer ID only     */
    H5PL_vol_key_t vol;           /* VOL: supports by-name        */
    H5PL_vfd_key_t vfd;           /* VFD: supports by-name        */
} H5PL_key_t;
\end{minted}

By contrast, VFDs use a discriminated union (\texttt{H5PLprivate.h}):

\begin{minted}{c}
typedef struct H5PL_vfd_key_t {
    H5FD_get_driver_kind_t kind;  /* BY_NAME or BY_VALUE          */
    union {
        H5FD_class_value_t value;
        const char        *name;
    } u;
} H5PL_vfd_key_t;
\end{minted}

When \texttt{H5FD\_register\_driver\_by\_name()} (\texttt{H5FDint.c}) needs to find a driver, it
sets \texttt{key.vfd.kind = H5FD\_GET\_DRIVER\_BY\_NAME}, passes the key to
\texttt{H5PL\_load()}, and the existing scan logic in \texttt{H5PL\_\_open()} loads the \texttt{.so},
calls \texttt{H5PLget\_plugin\_info()}, and invokes \texttt{H5FD\_check\_plugin\_load()} to do a
\texttt{strcmp} on the class name. The library already does an exhaustive scan of all
\texttt{.so} files in \texttt{HDF5\_PLUGIN\_PATH} for every plugin type---this is not a new
cost.

\verysubsection{Required changes}
Give filters the same discriminated union:

\begin{minted}{c}
typedef enum {
    H5Z_GET_FILTER_BY_ID,
    H5Z_GET_FILTER_BY_NAME
} H5Z_get_filter_kind_t;

typedef struct H5PL_filter_key_t {
    H5Z_get_filter_kind_t kind;
    union {
        H5Z_filter_t    id;
        const char     *name;
    } u;
} H5PL_filter_key_t;
\end{minted}

Replace the bare \texttt{int} in \texttt{H5PL\_key\_t}:

\begin{minted}{c}
typedef union H5PL_key_t {
    H5PL_filter_key_t filter;   /* NEW: by-name and by-id */
    H5PL_vol_key_t    vol;
    H5PL_vfd_key_t    vfd;
} H5PL_key_t;
\end{minted}

Add an \texttt{H5Z\_check\_plugin\_load()} callback (analogous to
\texttt{H5FD\_check\_plugin\_load} and \texttt{H5VL\_check\_plugin\_load}). In the
\texttt{H5PL\_TYPE\_FILTER} case within \texttt{H5PL\_\_open()}, this callback inspects the
class struct returned by \texttt{H5PLget\_plugin\_info()} and matches by either integer
ID or name depending on \texttt{key.filter.kind}. For name matching, comparison is
case-insensitive (ASCII folding).

Existing integer-ID search paths MUST be updated to use the new
\texttt{H5PL\_filter\_key\_t} with \texttt{kind = H5Z\_GET\_FILTER\_BY\_ID}. This is a mechanical
refactor with no behavioral change---the match logic is identical, just
accessed through the discriminated union. Note: this touches every existing
call site that currently writes \texttt{key.id = filter\_id} (in \texttt{H5Z.c},
\texttt{H5Pdcpl.c}, and any internal filter availability checks).

\verysubsection{Resolution procedure}
When \texttt{H5Pset\_filter\_by\_name} cannot resolve a name
from the in-memory registry:

\begin{enumerate}
\item Check whether the filter is already registered by name in the registry---if
  so, use the cached ID.
\item Set \texttt{key.filter.kind = H5Z\_GET\_FILTER\_BY\_NAME} and
  \texttt{key.filter.u.name = requested\_name}.
\item Call \texttt{H5PL\_load(H5PL\_TYPE\_FILTER, \&key)}.
\item The existing \texttt{H5PL\_\_find\_plugin\_in\_path\_table()} $\rightarrow$
  \texttt{H5PL\_\_open()} scan iterates directories in \texttt{HDF5\_PLUGIN\_PATH}. For each
  \texttt{.so}, it calls \texttt{H5PLget\_plugin\_type()} and \texttt{H5PLget\_plugin\_info()}.
\item \texttt{H5Z\_check\_plugin\_load()} compares the class struct's \texttt{name} field
  against the requested name (case-insensitive).
\item On match, the plugin is registered (both by ID and by name), and the
  filter ID is returned.
\item If no match is found after scanning all directories, \texttt{H5E\_NOFILTER} is
  returned.
\end{enumerate}

Once found, the name is cached in the registry. Subsequent calls resolve
instantly without scanning.

\verysubsection{v2 plugin behavior during name-based search}
For v2 plugins
(\texttt{H5Z\_class2\_t}), the \texttt{name} field is documented as ``Comment for debugging''
and may contain arbitrary strings. The by-name search MUST only match against
v3 plugins (\texttt{H5Z\_class3\_t}) where the \texttt{name} field is a canonical identifier.
v2 plugins are skipped during name-based search. They can still be discovered
by integer ID (the existing path) or pre-registered with \texttt{H5Zregister\_name}.

\verysubsection{Plugin path management}
Plugins discoverable by name are subject to the
same path configuration as all other plugins. The directories searched are
those configured via \texttt{HDF5\_PLUGIN\_PATH}, \texttt{H5PLappend()},
\texttt{H5PLprepend()}, and related \texttt{H5PL} path management APIs. No separate path
configuration is needed for name-based discovery.

\verysubsection{Windows platform considerations}
On Windows, plugins are DLLs rather than \texttt{.so} files, and
\texttt{HDF5\_PLUGIN\_PATH} uses the semicolon (\texttt{;}) as the directory separator
rather than the colon used on POSIX systems. The name-based scan iterates all
files in each plugin directory and attempts to load each candidate, which is
the same behavior as the existing integer-ID scan.


\subsection{Relationship to the HDF Group Filter Registry}

The HDF Group maintains an external registry of filter
IDs~\cite{hdf5registeredfilters}. Canonical names MUST be added to this
registry as a normative field before the first release that includes this
feature (see Q2 in \SectionRef{sec:open-questions}). This ensures global
uniqueness of filter names and provides plugin authors with an authoritative
source of truth for the \texttt{name} field in \texttt{H5Z\_class3\_t}.


\subsection{Registry Implementation Requirements}
\label{sec:registry-impl}

The name registry must support the following operations efficiently:

\begin{itemize}
\item \textbf{Forward lookup:} name $\rightarrow$ filter ID (used by
  \texttt{H5Pset\_filter\_by\_name}, \texttt{H5Zlookup\_by\_name}).
\item \textbf{Reverse lookup:} filter ID $\rightarrow$ canonical name (used by
  \texttt{H5Pget\_filter\_name\_by\_idx}).
\item \textbf{Reverse cleanup:} filter ID $\rightarrow$ all registered names (used by
  \texttt{H5Zunregister}, \SectionRef{sec:collision-policy}).
\item \textbf{Insertion-order tracking:} The introspection functions return ``the
  name that was first registered for that ID''
  (\SectionRef{sec:explicit-registration}). This requires either insertion-order
  metadata or a separate ID $\rightarrow$ first-name mapping. A simple
  name $\rightarrow$ ID hash map alone is not sufficient.
\end{itemize}

A bidirectional mapping (name $\rightarrow$ ID hash table plus ID $\rightarrow$ ordered name
list) satisfies all requirements. The exact data structure is an
implementation choice.

\end{document}
